\documentclass[11pt]{article}

\usepackage[margin=1in]{geometry}
\usepackage{amsmath,amssymb}
\usepackage{booktabs}
\usepackage{enumitem}
\usepackage[hidelinks]{hyperref}

\setlist[itemize]{leftmargin=1.4em}
\setlist[enumerate]{leftmargin=1.6em}

\title{SYSC 5108 Exam Study\\Assignment 2, Question 6}
\author{Jesse Levine}
\date{}

\begin{document}
\maketitle

\section*{Question}

Consider training a logistic regression model using a loss close to least squared
error (LSE). Assume \(K=2\), and define the differentiable loss through errors in
probabilities:
\[
    J(\beta)
    =
    \sum_{i=1}^{M}
    \left[
        I_1(y^{(i)})\left(1-p_1(z^{(i)})\right)^2
        +
        I_2(y^{(i)})\left(p_1(z^{(i)})\right)^2
    \right],
\]
where
\[
    z^{(i)} = \beta^T x^{(i)}, \qquad x_0=1.
\]
Find the gradient of the loss function and compare it with the MLE gradient
discussed in class.

\section*{Course and Textbook Cross-References}

\begin{itemize}
    \item \textbf{Chapter 3 slides, Logistic Regression Model:} for \(K=2\),
    logistic regression models the class probability using
    \(p_1(x)=\sigma(z)\), where \(z=\beta^T x\).
    \item \textbf{Chapter 3 slides, Alternative Notation:} for \(y\in\{0,1\}\),
    the negative log-likelihood loss has gradient
    \[
        \nabla_\beta J_{\text{MLE}}(\beta)
        =
        -\sum_i \left(y^{(i)}-\sigma(z^{(i)})\right)x^{(i)}.
    \]
    \item \textbf{Chapter 3 slides, Measuring Loss:} using negative
    log-likelihood is equivalent to cross-entropy for the model distribution.
    \item \textbf{Goodfellow, Bengio, and Courville, \emph{Deep Learning},
    Chapters 5 and 6:} logistic regression uses a sigmoid to map a linear score
    into \((0,1)\), and maximum likelihood becomes minimization of negative
    log-likelihood, equivalently cross-entropy. The textbook also notes that
    mean squared error with sigmoid outputs can suffer from saturation and slow
    learning.
\end{itemize}

\section*{Step-by-Step Solution}

\subsection*{1. Simplify the notation}

For one sample \(i\), define
\[
    p_i = p_1(z^{(i)}) = \sigma(z^{(i)}),
    \qquad
    z^{(i)} = \beta^T x^{(i)}.
\]
Also define the binary target
\[
    t_i = I_1(y^{(i)}).
\]
Since there are only two classes,
\[
    I_2(y^{(i)}) = 1 - I_1(y^{(i)}) = 1-t_i.
\]
Therefore, the sample loss can be rewritten as
\[
    J_i(\beta)
    =
    t_i(1-p_i)^2 + (1-t_i)p_i^2.
\]
This is exactly squared error on the predicted probability:
\[
    J_i(\beta) = (t_i-p_i)^2.
\]
Check:
\[
    \begin{cases}
    t_i=1: & J_i = (1-p_i)^2, \\
    t_i=0: & J_i = p_i^2.
    \end{cases}
\]

\subsection*{2. Differentiate the loss with respect to \(p_i\)}

\[
    J_i = (t_i-p_i)^2.
\]
Thus
\[
    \frac{\partial J_i}{\partial p_i}
    =
    2(p_i-t_i).
\]
Equivalently, using the original indicator notation:
\[
\begin{aligned}
    \frac{\partial J_i}{\partial p_i}
    &=
    -2I_1(y^{(i)})(1-p_i)+2I_2(y^{(i)})p_i.
\end{aligned}
\]

\subsection*{3. Differentiate the sigmoid with respect to \(\beta\)}

The sigmoid derivative is
\[
    \frac{d\sigma(z)}{dz}
    =
    \sigma(z)(1-\sigma(z)).
\]
Since \(z^{(i)}=\beta^T x^{(i)}\),
\[
    \nabla_\beta z^{(i)} = x^{(i)}.
\]
Therefore,
\[
    \nabla_\beta p_i
    =
    \frac{dp_i}{dz^{(i)}}\nabla_\beta z^{(i)}
    =
    p_i(1-p_i)x^{(i)}.
\]

\subsection*{4. Apply the chain rule}

For one sample,
\[
\begin{aligned}
    \nabla_\beta J_i(\beta)
    &=
    \frac{\partial J_i}{\partial p_i}\nabla_\beta p_i \\
    &=
    2(p_i-t_i)p_i(1-p_i)x^{(i)}.
\end{aligned}
\]
Summing over all \(M\) samples gives
\[
    \boxed{
    \nabla_\beta J(\beta)
    =
    2\sum_{i=1}^{M}
    \left(p_i-t_i\right)p_i(1-p_i)x^{(i)}
    }.
\]

Using \(t_i=I_1(y^{(i)})\) and \(p_i=p_1(z^{(i)})\), this is
\[
    \boxed{
    \nabla_\beta J(\beta)
    =
    2\sum_{i=1}^{M}
    \left(p_1(z^{(i)})-I_1(y^{(i)})\right)
    p_1(z^{(i)})\left(1-p_1(z^{(i)})\right)x^{(i)}
    }.
\]

In the assignment's original indicator notation:
\[
    \boxed{
    \nabla_\beta J(\beta)
    =
    2\sum_{i=1}^{M}
    \left[
        -I_1(y^{(i)})(1-p_i)
        +
        I_2(y^{(i)})p_i
    \right]
    p_i(1-p_i)x^{(i)}
    }.
\]

\section*{Comparison With MLE}

For binary logistic regression, the MLE or negative-log-likelihood loss is
\[
    J_{\text{MLE}}(\beta)
    =
    -\sum_{i=1}^{M}
    \left[
        t_i\log(p_i)
        +
        (1-t_i)\log(1-p_i)
    \right].
\]
Its gradient is
\[
    \boxed{
    \nabla_\beta J_{\text{MLE}}(\beta)
    =
    \sum_{i=1}^{M}(p_i-t_i)x^{(i)}
    =
    -\sum_{i=1}^{M}(t_i-p_i)x^{(i)}
    }.
\]

Now compare the two gradients:
\[
    \nabla_\beta J_{\text{LSE-prob}}(\beta)
    =
    2\sum_i (p_i-t_i)\underbrace{p_i(1-p_i)}_{\text{extra sigmoid derivative}}x^{(i)}
\]
while
\[
    \nabla_\beta J_{\text{MLE}}(\beta)
    =
    \sum_i (p_i-t_i)x^{(i)}.
\]

\section*{Interpretation}

\begin{itemize}
    \item The LSE-on-probabilities gradient has an extra factor
    \(p_i(1-p_i)\), and \(p_i(1-p_i)\le 0.25\).
    \item When the sigmoid saturates, \(p_i\approx 0\) or \(p_i\approx 1\), this
    extra factor becomes very small.
    \item This is a serious problem when the model is \emph{confident but wrong}.
    For example, if \(t_i=1\) but \(p_i\approx 0\), then
    \[
        \nabla_\beta J_{\text{MLE},i}\approx -x^{(i)},
    \]
    which gives a strong corrective update, but
    \[
        \nabla_\beta J_{\text{LSE-prob},i}
        \approx
        -2p_i(1-p_i)x^{(i)}
        \approx 0,
    \]
    which can make learning very slow.
    \item This is why MLE/cross-entropy is preferred for logistic regression and
    sigmoid output units: the log-likelihood cancels the sigmoid derivative in a
    way that avoids this extra saturation factor.
\end{itemize}

\section*{Final Answer}

\[
    \boxed{
    \nabla_\beta J(\beta)
    =
    2\sum_{i=1}^{M}
    \left(p_1(z^{(i)})-I_1(y^{(i)})\right)
    p_1(z^{(i)})\left(1-p_1(z^{(i)})\right)x^{(i)}
    }
\]

Compared with MLE,
\[
    \boxed{
    \nabla_\beta J_{\text{MLE}}(\beta)
    =
    \sum_{i=1}^{M}
    \left(p_1(z^{(i)})-I_1(y^{(i)})\right)x^{(i)}
    }.
\]
The LSE-style probability loss therefore has the same prediction-error term
\((p_i-t_i)\), but it also has the extra sigmoid-derivative factor
\(p_i(1-p_i)\), which can cause vanishing gradients for saturated sigmoid outputs.

\section*{Exam Note}

If the loss is written with a leading factor \(\tfrac12\), then the factor \(2\)
in the LSE gradient disappears:
\[
    \frac{1}{2}(t_i-p_i)^2
    \quad\Rightarrow\quad
    \nabla_\beta J_i
    =
    (p_i-t_i)p_i(1-p_i)x^{(i)}.
\]
The key comparison with MLE does not change: the LSE probability loss still has
the extra \(p_i(1-p_i)\) factor.

\end{document}
